Indirect Time-of-Flight (iToF) sensors measure depth by correlating the
phase shift of amplitude-modulated illumination, providing metric depth at
lower cost and power than mechanical LiDAR.  Modern iToF sensors are
ubiquitous in smartphones (e.g., Apple LiDAR with $\leq$576 points) and
emerging consumer devices (as few as 64 points at 200\,mW).  This extreme
sparsity---potentially $8\times8$ depth points per frame---poses a unique
challenge for Gaussian Splatting SLAM, which traditionally requires dense
depth for initialization and supervision.

\subsection{ToF-Splatting}

ToF-Splatting~\cite{conti2025_tofsplatting} (CVPR 2025) is the first 3DGS-based
SLAM pipeline tailored for very sparse ToF input.  Its architecture consists
of three modules:

\begin{enumerate}
    \item \textbf{Tracking Frontend}: A frame-to-model approach that
    minimizes a combined photometric ($\mathcal{L}_{\text{rgb}}$) and
    geometric ($\mathcal{L}_{\text{depth}}$) loss:
    \begin{equation}
        \mathcal{L}_{\text{track}} = \lambda_{\text{track}}
        \mathcal{L}_{\text{rgb}} + (1-\lambda_{\text{track}})
        \mathcal{L}_{\text{depth}}
    \end{equation}
    where $\mathcal{L}_{\text{depth}}$ uses the dense depth predicted by the
    multi-frame integration module (not the raw sparse ToF depth).

    \item \textbf{Multi-Frame Integration (MFI)}: Extends the Depth on Demand
    (DoD) framework~\cite{conti2024_dod} to fuse three complementary cues:
    (a)~extremely sparse ToF depth, (b)~monocular color cues via a Depth
    Anything v2~\cite{yang2024depthanythingv2} encoder, and
    (c)~multi-view geometry from a keyframe buffer.  An outlier rejection
    step compares ToF measurements against a no-ToF depth prediction and
    discards measurements with $\ell_1$ error above a quantile threshold,
    producing filtered sparse depth $\tilde{H}_k$ before final DoD
    prediction.

    \item \textbf{Mapping Backend}: Initializes new Gaussians where rendered
    opacity is low (uncertainty-driven).  The mapping loss combines
    photometric, DSSIM, depth, normal, and isotropy regularization terms:
    \begin{equation}
        \mathcal{L}_{\text{map}} = \lambda_{\text{map}}[\lambda_{\text{visual}}
        \mathcal{L}_{\text{rgb}} + (1-\lambda_{\text{visual}})
        \mathcal{L}_{\text{dssim}}] + (1-\lambda_{\text{map}})
        \mathcal{L}_{\text{depth}} + \lambda_{\text{normals}}
        \mathcal{L}_{\text{normals}} + \lambda_{\text{iso}}
        \mathcal{L}_{\text{iso}}
    \end{equation}
    The isotropy loss $\mathcal{L}_{\text{iso}}$ penalizes elongated
    Gaussians, and a keyframe novelty factor $\nu_k$ based on rendered
    opacity drives keyframe selection.
\end{enumerate}

\textbf{Key results}: On the ZJUL5 dataset (real sparse ToF) and Replica,
ToF-Splatting achieves state-of-the-art tracking and mapping, outperforming
prior sparse-ToF SLAM (ToF-SLAM~\cite{jiang2022quantizedtof}) and achieving
comparable or better results than RGB-D methods even with $\leq$1\% of the
depth points.  On TUM RGB-D, it matches full-density RGB-D Gaussian SLAM.

\textbf{Code}: Project page at
\url{https://andreaconti.github.io/projects/tof_splatting}.

\subsection{iToF-RGB Depth Fusion (Non-SLAM)}

While not SLAM systems per se, several works address the fundamental
iToF+RGB fusion problem that underpins this category:

\textbf{Fang et al.}~\cite{fang2025_iToFRGB} (ICPR 2024) propose a
weakly-supervised framework fusing narrow-FoV iToF depth with wide-FoV RGB
monocular depth.  The learning objective is inter-camera geometric
consistency, supplemented by structure distillation from the passive
sensing modality.  This is the smartphone-targeted counterpart to
ToF-Splatting.

\textbf{Du et al.}~\cite{du2025_iToFRGB} (arXiv 2025) present an iToF-RGB
fusion framework addressing iToF's low spatial resolution, limited FoV, and
structural distortion.  A precise geometric calibration module reprojects
narrow-FoV iToF depth onto the wide-FoV RGB coordinate system, and a
dual-encoder fusion network jointly extracts complementary features guided
by monocular depth priors.

\subsection{ToF-IP: Inertial + ToF for Motion Capture}

ToF-IP~\cite{yao2025_tofip} (NeurIPS 2025) integrates ToF sensors with
sparse IMUs for human motion capture---a different application domain but
demonstrating the complementary nature of ToF (direct distance measurement)
and inertial sensing (indirect acceleration integration).  A Transformer
encoder with Node-Centric Data Integration and Dynamic Spatial Positional
Encoding effectively fuses the two modalities.

\subsection{Summary and Gap Analysis}

The iToF + RGB + IMU Gaussian SLAM category is currently the least explored
of the three, with ToF-Splatting being the sole dedicated 3DGS-SLAM system.
The key gaps are:

\begin{itemize}
    \item \textbf{No IMU integration}: ToF-Splatting uses only ToF + RGB,
    not IMU.  An iToF + RGB + IMU Gaussian SLAM system would combine the
    metric scale from iToF, motion continuity from IMU, and appearance from
    RGB---a natural and unexplored configuration for mobile/AR devices.
    \item \textbf{Outdoor extension}: iToF sensors suffer from ambient light
    interference outdoors.  Extending to outdoor scenes requires robust
    multi-frame integration under varying SNR.
    \item \textbf{Embedded deployment}: iToF sensors target mobile devices,
    but current 3DGS optimization is GPU-intensive.  Lightweight Gaussian
    representations (e.g., LoCoSplat~\cite{wang2026_locosplat}) may enable
    on-device real-time mapping.
\end{itemize}

\begin{table}[t]
\centering
\caption{iToF + RGB + IMU Gaussian Splatting and related systems.}
\label{tab:itof_summary}
\scriptsize
\begin{tabular}{lcccc}
\toprule
\textbf{System} & \textbf{Year} & \textbf{Venue} & \textbf{3DGS SLAM} & \textbf{IMU} \\
\midrule
ToF-Splatting~\cite{conti2025_tofsplatting} & 2025 & CVPR & \checkmark & -- \\
ToF-SLAM~\cite{jiang2022quantizedtof} & 2022 & CVPR-W & -- (NeRF) & -- \\
DoD~\cite{conti2024_dod} & 2024 & ECCV & -- (depth only) & -- \\
iToF-RGB Fusion~\cite{fang2025_iToFRGB} & 2024 & ICPR & -- & -- \\
iToF-RGB Fusion~\cite{du2025_iToFRGB} & 2025 & arXiv & -- & -- \\
ToF-IP~\cite{yao2025_tofip} & 2025 & NeurIPS & -- (MoCap) & \checkmark \\
\bottomrule
\end{tabular}
\end{table}
